\section*{Chapitre \ref{ch:g12:subduction-magmatism} --- Le magmatisme des zones de subduction}

\begin{solution}{exo:g12:subduction-magmatism:1}
(a) un \omterm{def:g12:subduction-magmatism:arc}{arc volcanique}~; (b) la \omterm{def:g12:subduction-magmatism:hydrous-melting}{fusion hydratée} (fusion par apport d'eau)~;
(c) une \omterm{def:g12:subduction-magmatism:granodiorite}{granodiorite}~; (d) un \omterm{def:g12:subduction-magmatism:batholith}{batholite}~; (e) la \omterm{def:g12:subduction-magmatism:continental-growth}{croissance continentale}.
\end{solution}

\begin{solution}{exo:g12:subduction-magmatism:2}
(1) Au moins $24/0.25 = \qty{96}{km/h}$, soit $\num{96000}/3600 \approx
\qty{27}{m/s}$. (2) Le souffle était au moins $480/96 = 5$ fois plus
rapide.
\end{solution}

\begin{solution}{exo:g12:subduction-magmatism:3}
(1) $838/1214 \approx 0.690$, soit environ \qty{69}{\percent}. (2) Du côté
de la plaque chevauchante, celle qui ne plonge pas. (3) Entre la fosse et
le front, le \omterm{def:g11:subduction-zones:slab}{panneau plongeant} n'est pas encore assez profond~: il n'a pas
atteint la profondeur, une centaine de kilomètres, où son eau est libérée
dans un manteau assez chaud pour fondre une fois hydraté.
\end{solution}

\begin{solution}{exo:g12:subduction-magmatism:4}
(1) $45/(45 + 15) = 0.75$~: \qty{75}{\percent} du feldspath est du
plagioclase. (2) Oui~: elle est à gros grains, contient plus de
\qty{20}{\percent} de quartz, et sa fraction de plagioclase est comprise
entre 65 et \qty{90}{\percent}.
\end{solution}

\begin{solution}{exo:g12:subduction-magmatism:5}
Les magmas d'arc sont plus riches en silice, donc plus visqueux~; et ils
transportent l'eau qui les a produits, donc ils sont riches en gaz. Une
forte viscosité et une forte teneur en gaz favorisent toutes deux les
éruptions explosives.
\end{solution}

\begin{solution}{exo:g12:subduction-magmatism:6}
(1) $H = 150 \times \tan\ang{35} \approx 150 \times 0.700 \approx
\qty{105}{km}$. (2) C'est la moyenne mondiale.
\end{solution}

\begin{solution}{exo:g12:subduction-magmatism:7}
$d = H/\tan\theta$~: $105/\tan\ang{20} \approx 105/0.364 \approx
\qty{288}{km}$ à \ang{20}, et $105/\tan\ang{45} = \qty{105}{km}$ à
\ang{45}. Plus le \omterm{def:g11:subduction-zones:slab}{panneau} est incliné doucement, plus l'arc se trouve loin
de la fosse.
\end{solution}

\begin{solution}{exo:g12:subduction-magmatism:8}
(1) $T = -5.104 \times 3.16^2 + 132.9 \times 3.16 + 1120.7 \approx -51.0 +
420.0 + 1120.7 \approx \qty{1490}{\celsius}$. (2) D'environ $1490 - 1000 =
\qty{490}{\celsius}$. (3) À sec, non~: \qty{1300}{\celsius} est en dessous
de \qty{1490}{\celsius}. En présence d'eau, oui~: c'est au-dessus de
\qty{1000}{\celsius}.
\end{solution}

\begin{solution}{exo:g12:subduction-magmatism:9}
(1) $2.5/2.7 \approx 0.93$ an, soit environ 340 jours. (2) L'avalanche a
déplacé des roches anciennes qui existaient déjà~; les arcs fabriquent de
la roche nouvelle à partir du manteau. On ne compare que les volumes.
\end{solution}

\begin{solution}{exo:g12:subduction-magmatism:10}
(1) $2.7 \times \num{1000000} = \qty{2.7e6}{km^3}$. (2) $\num{9e9}/2.7
\approx \num{3.3e9}$ ans, soit 3.3~milliards d'années. (3) C'est moins que
l'âge de la Terre~: les arcs auraient donc pu fabriquer toute la \omterm{def:g12:continental-crust:continental-crust}{croûte
continentale}~; mais la \omterm{def:g11:subduction-zones:subduction}{subduction} en reprend aussi, si bien qu'en réalité
une grande partie de la croûte a été fabriquée et détruite plus d'une fois.
\end{solution}

\begin{solution}{exo:g12:subduction-magmatism:11}
(1) Des séismes superficiels se produisent aussi dans la plaque
chevauchante et sur la faille qui sépare les deux plaques~; seuls les
séismes profonds se trouvent dans le \omterm{def:g11:subduction-zones:slab}{panneau plongeant} seul.
(2) $116 - 80 = \qty{36}{km}$ sur \qty{86.4}{km}~: $\tan\theta = 36/86.4
\approx 0.417$, $\theta \approx \ang{22.6}$.
\end{solution}

\begin{solution}{exo:g12:subduction-magmatism:12}
Dans le triangle rectangle de la proposition, le \omterm{def:g11:subduction-zones:slab}{panneau} est l'hypoténuse
$L$ et $d$ le côté adjacent à l'angle $\theta$~: $\cos\theta = d/L$ et $L =
d/\cos\theta$. Pour le Japon, $L = 270/\cos\ang{20} \approx 270/0.940 \approx
\qty{287}{km}$. À \qty{83}{mm} par an~: $\num{2.87e8}/83 \approx \num{3.5e6}$
ans, soit environ 3.5 millions d'années.
\end{solution}

\begin{solution}{exo:g12:subduction-magmatism:13}
(1) $V = \pi r^2 h = \pi \times 0.5^2 \times 8 \approx \qty{6.3}{km^3}$.
(2) Environ $6.3/2.5 \approx 2.5$ fois le volume de l'avalanche.
(3) $P = \rho g z$~: à \qty{6}{km}, $2830 \times 9.81 \times 6000 \approx
\qty{1.67e8}{Pa} = \qty{0.17}{GPa}$~; à \qty{14}{km}, environ
\qty{3.89e8}{Pa} = \qty{0.39}{GPa}.
\end{solution}

\begin{solution}{exo:g12:subduction-magmatism:14}
(1) $2.7 - 3.2 = \qty{-0.5}{km^3}$ par an, soit une perte de
\qty{5e5}{km^3} par million d'années. (2) $500 \times \num{5e5} = \qty{2.5e8}{km^3}$,
soit $\num{2.5e8}/\num{9e9} \approx \qty{2.8}{\percent}$ de la croûte.
(3) Chaque rythme est une estimation, incertaine de bien plus que l'écart
de \qty{0.5}{km^3} par an qui les sépare~: aux erreurs près, les deux
s'équilibrent.
\end{solution}

\begin{solution}{exo:g12:subduction-magmatism:15}
(1) $\num{20000} \times 10 = \qty{2e5}{km^3}$. (2) $\num{2e5}/1 =
\num{200000}$ ans pour tous les arcs continentaux du monde réunis. Un seul
arc ne reçoit qu'une petite part de ce magma~: un \omterm{def:g12:subduction-magmatism:batholith}{batholite} met donc bien
plus longtemps à se former~; il se construit \omterm{def:g11:igneous-rocks:intrusion}{pluton} après \omterm{def:g11:igneous-rocks:intrusion}{pluton}, sur des
millions d'années. (3) Au moins les \qty{5}{km} de roches situés au-dessus
de son sommet.
\end{solution}

\begin{solution}{pb:g12:subduction-magmatism:1}
\textbf{1.} $111.2 \times \cos\ang{39} \approx 111.2 \times 0.777 \approx
\qty{86.4}{km}$.

\textbf{2.} $(144.1 - 141.0) \times 86.4 \approx 3.1 \times 86.4 \approx
\qty{268}{km}$, environ \qty{270}{km}.

\textbf{3.} $(141.0 - 139.7) \times 86.4 \approx \qty{112}{km}$.

\textbf{4.} Sur la plaque chevauchante, celle qui ne plonge pas, à l'ouest
de la fosse~: la plaque Pacifique plonge sous elle.

\textbf{5.} À cet endroit, le \omterm{def:g11:subduction-zones:slab}{panneau plongeant} n'est pas encore assez
profond pour libérer son eau dans un manteau assez chaud pour fondre.

\textbf{6.} $116 - 80 = \qty{36}{km}$.

\textbf{7.} Un degré de longitude, \qty{86.4}{km}.

\textbf{8.} $\tan\theta = 36/86.4 \approx 0.417$, donc $\theta \approx
\ang{22.6}$.

\textbf{9.} La droite remonte de \qty{36}{km} tous les \qty{86.4}{km} vers
l'est~: elle atteint donc la surface $98/0.417 \approx \qty{235}{km}$ à
l'est de l'Iwate, à $141.0 + 235/86.4 \approx \ang{143.7}$~E, soit environ $(144.1 - 143.7) \times 86.4
\approx \qty{35}{km}$ à l'ouest de la fosse.

\textbf{10.} Le \omterm{def:g11:subduction-zones:slab}{panneau} n'est pas un plan~: il entre dans la fosse presque
à l'horizontale et s'incurve de plus en plus vers le bas à mesure qu'il
plonge.

\textbf{11.} $268 \times \tan\ang{20} \approx 268 \times 0.364 \approx
\qty{98}{km}$.

\textbf{12.} Elles concordent~: \qty{98}{km} dans les deux cas.

\textbf{13.} Les deux se situent dans l'intervalle mondial, un peu moins
profond que la moyenne de \qty{105}{km}.

\textbf{14.} Le Megata se trouve \qty{112}{km} à l'ouest de l'Iwate~: le
\omterm{def:g11:subduction-zones:slab}{panneau} est donc sous lui à environ
$98 + 0.417 \times 112 \approx \qty{145}{km}$ de profondeur.

\textbf{15.} Dans le manteau, environ $98 - 35 = \qty{63}{km}$ sous la base
de la croûte.

\textbf{16.} Elle est entrée près de la dorsale où la croûte s'est formée, à
partir de l'eau de mer~; une grande partie en est retenue dans la
serpentine.

\textbf{17.} $T = -5.104 \times 3.16^2 + 132.9 \times 3.16 + 1120.7 \approx
\qty{1490}{\celsius}$.

\textbf{18.} De $1490 - 1000 = 490$ à $1490 - 950 = \qty{540}{\celsius}$~:
environ \qty{500}{\celsius}.

\textbf{19.} Le \omterm{def:g11:subduction-zones:slab}{panneau} est de plusieurs centaines de degrés plus froid que
le manteau qui l'entoure, trop froid pour fondre même en présence
d'eau~: c'est le coin de manteau chaud situé au-dessus, une fois hydraté,
qui fond.

\textbf{20.} Le \omterm{def:g11:subduction-zones:slab}{panneau plongeant} se trouve à environ \textbf{\qty{100}{km}
sous le front volcanique} du nord-est du Japon (\qty{98}{km} par les deux
méthodes).
\end{solution}
